%=======================================================
\chapter{Aspas e outros símbolos tipográficos}

\textsc{Pacote:} \textbf{textcmds} (Acesso: \url{https://ctan.math.illinois.edu/macros/latex/contrib/amsrefs/textcmds.pdf})

\textsc{O que faz?}: Permite incluir aspas duplas e simples, e outros símbolos tipográficos, com menos esforço na digitação. Ver explicação a seguir.

Por padrão, o \LaTeX\ produz aspas das seguintes formas:

\begin{enumerate}
	\item SHIFT + acento grave 2x + SHIFT + acento grave 2x + palavra ou expressão entre aspas + aspas simples 2x.
	\begin{itemize}
		\item 'um texto qualquer' entre aspas simples
	\end{itemize}
	\item SHIFT + sinal de acento grave + palavra ou expressão + aspas simples.
	\begin{itemize}
		\item "um texto qualquer" entre aspas duplas
	\end{itemize}
\end{enumerate}

No primeiro exemplo, o 1º sinal das aspas simples não está correto, pois deveria estar voltado para direita.

No segundo exemplo, não acrescenta o espaço entre as últimas aspas duplas e a próxima palavra.

Para resolver isso e utilizar ambas as aspas corretamente, use os comandos:
\ \\

\begin{codex}{Códigos para aspas com o pacote textcmds}
Estas são \qq{aspas duplas}.
Estas são \q{aspas simples}.
Estas são \qq{aspas duplas com \q{aspas simples} ao mesmo tempo}
\end{codex}
\ \\

\textsc{Resultado}:

Estas são \qq{aspas duplas}; Estas são \q{aspas simples} e estas são \qq{aspas duplas com \q{aspas simples} ao mesmo tempo}.
